\documentclass{exam-zh}
\usepackage{edu-practice}

\examsetup{
  question/show-answer = false,
  fillin/show-answer = false,
  solution/show-solution = hide,
}

\begin{document}

\section*{2026 届一模选择填空压轴精选}

% ---------- 第 1 题 (嘉定·选择6) ----------
\needspace{10\baselineskip}

\begin{question}[points=4]
\noindent
  \begin{minipage}[t]{\dimexpr\linewidth-65mm-6mm\relax}
    \vspace{0pt}
    \hfill{\color{edu-blue!60}\small\textit{嘉定·选择6}}\par
    如图2，正方形 $DEFG$ 的边 $EF$ 在 $\triangle ABC$ 的边 $BC$ 上，顶点 $D$、$G$ 分别在边 $AB$、$AC$ 上，已知 $BC=AH$，选择不正确的结论。
    \par
    \begin{choices}[columns=1]
      \item $\triangle ABC\sim\triangle ADG$
      \item $\dfrac{AD}{AB}=\dfrac{AP}{AH}$
      \item $S_{\triangle ADG}=\dfrac25S_{\triangle ABC}$
      \item $AP=PH$
    \end{choices}
  \end{minipage}\hfill
  \begin{diagramcoltikz}{65mm}{}
  \includegraphics[width=\linewidth]{\detokenize{../2026-07-24-上海初三试卷原题库/staging/2026-JIADING-TERM/items/Q006/assets/prompt-figure.png}}
  \end{diagramcoltikz}
  \par\medskip
  \answerarea[40mm]
\end{question}

% ---------- 第 2 题 (嘉定·填空18) ----------
\needspace{10\baselineskip}

\begin{question}[points=4]
\noindent
  \begin{minipage}[t]{\dimexpr\linewidth-65mm-6mm\relax}
    \vspace{0pt}
    \hfill{\color{edu-blue!60}\small\textit{嘉定·填空18}}\par
    如图8，$Rt\triangle ABC$ 中，已知 $\angle C=90^\circ$，$AC=2$，$BC=4$，将 $\triangle ABC$ 绕点 $A$ 逆时针旋转 $\alpha\ (0^\circ<\alpha<90^\circ)$ 得到 $\triangle AB'C'$，点 $B$、$C$ 的对应点分别为 $B'$、$C'$，联结 $B'C$。如果直线 $B'C$ 垂直于 $AB$，联结 $B'B$，那么此时 $\triangle AB'B$ 的面积是 $\fillin$。
  \end{minipage}\hfill
  \begin{diagramcoltikz}{65mm}{}
  \includegraphics[width=\linewidth]{\detokenize{../2026-07-24-上海初三试卷原题库/staging/2026-JIADING-TERM/items/Q018/assets/prompt-figure.png}}
  \end{diagramcoltikz}
  \par\medskip
  \answerarea[48mm]
\end{question}

\newpage

% ---------- 第 3 题 (金山·选择5) ----------
\needspace{10\baselineskip}

\begin{question}[points=4]
    \hfill{\color{edu-blue!60}\small\textit{金山·选择5}}\par
    在等边 $\triangle ABC$ 中，点 $E$、$F$ 分别在边 $AB$、$AC$ 上，将 $\triangle ABC$ 沿 $EF$ 折叠，使得点 $A$ 与 $\triangle ABC$ 的重心 $G$ 重合，$EF$ 与 $AG$ 交于点 $O$，延长 $AG$ 交 $BC$ 于点 $D$，那么 $EF:OD$ 的值为（　）。
    \par
    \begin{choices}[columns=1]
      \item $\frac{\sqrt2}{2}$
      \item $\frac{\sqrt3}{2}$
      \item $\frac{\sqrt3}{3}$
      \item $1$
    \end{choices}
  \par\medskip
  \answerarea[40mm]
\end{question}

% ---------- 第 4 题 (金山·填空15) ----------
\needspace{10\baselineskip}

\begin{question}[points=4]
    \hfill{\color{edu-blue!60}\small\textit{金山·填空15}}\par
    在矩形 $ABCD$ 中，过点 $C$ 作 $CE\perp BD$，垂足为 $E$，以 $CE$ 为斜边作直角三角形 $CEF$，$AF$ 与 $BD$ 交于点 $P$。如果 $\dfrac{AP}{AF}=k$，那么 $k$ 的取值范围是 $\fillin$。
  \par\medskip
  \answerarea[48mm]
\end{question}

\newpage

% ---------- 第 5 题 (普陀·选择6) ----------
\needspace{10\baselineskip}

\begin{question}[points=4]
\noindent
  \begin{minipage}[t]{\dimexpr\linewidth-65mm-6mm\relax}
    \vspace{0pt}
    \hfill{\color{edu-blue!60}\small\textit{普陀·选择6}}\par
    如图，在等边三角形 $ABC$ 中，点 $D$ 在边 $BC$ 上，点 $D$ 关于直线 $AB、AC$ 的对称点分别是点 $P、Q$，那么关于线段 $PQ$ 与线段 $AD$ 之间的数量关系正确的是
    \par
    \begin{choices}[columns=1]
      \item $PQ=\dfrac{2\sqrt3}{3}AD$
      \item $PQ=\sqrt3AD$
      \item $PQ=2AD$
      \item $PQ=\dfrac{\sqrt3}{2}AD$
    \end{choices}
  \end{minipage}\hfill
  \begin{diagramcoltikz}{65mm}{}
  \includegraphics[width=\linewidth]{\detokenize{../2026-07-24-上海初三试卷原题库/staging/2026-PUTUO-TERM/items/Q006/assets/manual-prompt-1-cb991a967c4c8166.png}}
  \end{diagramcoltikz}
  \par\medskip
  \answerarea[40mm]
\end{question}

% ---------- 第 6 题 (普陀·填空18) ----------
\needspace{10\baselineskip}

\begin{question}[points=4]
\noindent
  \begin{minipage}[t]{\dimexpr\linewidth-65mm-6mm\relax}
    \vspace{0pt}
    \hfill{\color{edu-blue!60}\small\textit{普陀·填空18}}\par
    如图，$\triangle ABC$ 中，$\angle BAC=90^\circ$，$AC=3$，$AB=4$，点 $D$ 在边 $AB$ 上，将 $\triangle ACD$ 沿着 $CD$ 翻折得 $\triangle A'CD$，其中点 $A$ 与点 $A'$ 对应，联结 $A'B$。如果 $\angle ACD=\angle A'BA$，那么 $AD=\fillin$。
  \end{minipage}\hfill
  \begin{diagramcoltikz}{65mm}{}
  \includegraphics[width=\linewidth]{\detokenize{../2026-07-24-上海初三试卷原题库/staging/2026-PUTUO-TERM/items/Q018/assets/manual-prompt-1-29e7b26a7ba7c171.png}}
  \end{diagramcoltikz}
  \par\medskip
  \answerarea[48mm]
\end{question}

\newpage

% ---------- 第 7 题 (闵行·选择6) ----------
\needspace{10\baselineskip}

\begin{question}[points=4]
\noindent
  \begin{minipage}[t]{\dimexpr\linewidth-65mm-6mm\relax}
    \vspace{0pt}
    \hfill{\color{edu-blue!60}\small\textit{闵行·选择6}}\par
    如图，抛物线 $C_1:y=x^2-6x$ 向右平移得到 $C_2,C_3$，三条抛物线与直线 $y=m$ 分别交于两点。六个交点横坐标之和是
    \par
    \begin{choices}[columns=1]
      \item $6$
      \item $18$
      \item $30$
      \item $54$
    \end{choices}
  \end{minipage}\hfill
  \begin{diagramcoltikz}{65mm}{}
  \includegraphics[width=\linewidth]{\detokenize{../2026-07-24-上海初三试卷原题库/staging/2026-MINHANG-TERM/items/Q006/assets/prompt-01.png}}
  \end{diagramcoltikz}
  \par\medskip
  \answerarea[40mm]
\end{question}

% ---------- 第 8 题 (闵行·填空18) ----------
\needspace{10\baselineskip}

\begin{question}[points=4]
\noindent
  \begin{minipage}[t]{\dimexpr\linewidth-65mm-6mm\relax}
    \vspace{0pt}
    \hfill{\color{edu-blue!60}\small\textit{闵行·填空18}}\par
    如图，在矩形 $ABCD$ 中，$E$ 是 $BC$ 中点，过 $E$ 作 $EF\parallel BD$ 交 $CD$ 于 $F$，将 $\triangle CEF$ 沿 $EF$ 翻折到 $\triangle GEF$，若 $G$ 在 $AE$ 上，那么 $BG:AE=$\fillin。
  \end{minipage}\hfill
  \begin{diagramcoltikz}{65mm}{}
  \includegraphics[width=\linewidth]{\detokenize{../2026-07-24-上海初三试卷原题库/staging/2026-MINHANG-TERM/items/Q018/assets/prompt-01.png}}
  \end{diagramcoltikz}
  \par\medskip
  \answerarea[48mm]
\end{question}

\newpage

% ---------- 第 9 题 (松江·选择6) ----------
\needspace{10\baselineskip}

\begin{question}[points=4]
    \hfill{\color{edu-blue!60}\small\textit{松江·选择6}}\par
    命题①腰和腰上的高对应成比例的两个等腰三角形相似；命题②底边和腰上的高对应成比例的两个等腰三角形相似。判断正确的是
    \par
    \begin{choices}[columns=1]
      \item ①和②都是真命题
      \item ①和②都是假命题
      \item ①是真命题，②是假命题
      \item ①是假命题，②是真命题
    \end{choices}
  \par\medskip
  \answerarea[40mm]
\end{question}

% ---------- 第 10 题 (松江·填空18) ----------
\needspace{10\baselineskip}

\begin{question}[points=4]
    \hfill{\color{edu-blue!60}\small\textit{松江·填空18}}\par
    在 $Rt\triangle ABC$ 中，$\angle ACB=90^\circ$，点 $P,O$ 分别在 $AB,BC$ 上；若 $\triangle ACO\sim\triangle BPO$ 且 $\triangle APO$ 是等腰三角形，则相应角的正切值为 $\fillin$。
  \par\medskip
  \answerarea[48mm]
\end{question}

\newpage

% ---------- 第 11 题 (静安·选择6) ----------
\needspace{10\baselineskip}

\begin{question}[points=4]
\noindent
  \begin{minipage}[t]{\dimexpr\linewidth-58mm-6mm\relax}
    \vspace{0pt}
    \hfill{\color{edu-blue!60}\small\textit{静安·选择6}}\par
    如图，在矩形 $ABCD$ 中，$AB=6$，$BC=8$。分别以 $A、C$ 为圆心的两圆相切，且点 $D$ 在圆 $C$ 内、点 $B$ 在圆 $C$ 外。圆 $A$ 的半径 $r$ 的取值范围是
    \par
    \begin{choices}[columns=1]
      \item $6<r<8$ 或 $10<r<18$
      \item $2<r<4$ 或 $10<r<18$
      \item $2<r<4$ 或 $16<r<18$
      \item $6<r<8$ 或 $16<r<18$
    \end{choices}
  \end{minipage}\hfill
  \begin{diagramcoltikz}{58mm}{}
  \includegraphics[width=\linewidth]{\detokenize{../2026-07-24-上海初三试卷原题库/staging/2026-JINGAN-TERM/items/Q006/assets/prompt-01.png}}
  \end{diagramcoltikz}
  \par\medskip
  \answerarea[40mm]
\end{question}

% ---------- 第 12 题 (静安·填空18) ----------
\needspace{10\baselineskip}

\begin{question}[points=4]
\noindent
  \begin{minipage}[t]{\dimexpr\linewidth-58mm-6mm\relax}
    \vspace{0pt}
    \hfill{\color{edu-blue!60}\small\textit{静安·填空18}}\par
    如图，在等边 $\triangle ABC$ 中，点 $D、E$ 分别在边 $BC、AC$ 上，$AE=CD$，$AD$ 交 $BE$ 于点 $P$，$BQ\perp AD$ 于点 $Q$。如果 $PQ=3PE$，且 $PE=a$，那么等边 $\triangle ABC$ 的边长为\fillin。
  \end{minipage}\hfill
  \begin{diagramcoltikz}{58mm}{}
  \includegraphics[width=\linewidth]{\detokenize{../2026-07-24-上海初三试卷原题库/staging/2026-JINGAN-TERM/items/Q018/assets/prompt-01.png}}
  \end{diagramcoltikz}
  \par\medskip
  \answerarea[48mm]
\end{question}

\newpage

% ---------- 第 13 题 (崇明·选择6) ----------
\needspace{10\baselineskip}

\begin{question}[points=4]
    \hfill{\color{edu-blue!60}\small\textit{崇明·选择6}}\par
    在 $\triangle ABC$ 和 $\triangle DEF$ 中，$AB=AC$，$DE=DF$，根据下列条件，一定能判断 $\triangle ABC$ 和 $\triangle DEF$ 相似的是
    \par
    \begin{choices}[columns=1]
      \item $\frac{AB}{DE}=\frac{AC}{DF}$
      \item $\frac{AB}{DE}=\frac{AC}{EF}$
      \item $\angle B=\angle D$
      \item $\angle C=\angle E$
    \end{choices}
  \par\medskip
  \answerarea[40mm]
\end{question}

% ---------- 第 14 题 (崇明·填空18) ----------
\needspace{10\baselineskip}

\begin{question}[points=4]
\noindent
  \begin{minipage}[t]{\dimexpr\linewidth-65mm-6mm\relax}
    \vspace{0pt}
    \hfill{\color{edu-blue!60}\small\textit{崇明·填空18}}\par
    定义：当一个三角形有两个内角的差为 $90^\circ$ 时，这个三角形叫做“差直角三角形”。如图，在 $\triangle ABC$ 中，$AC=BC=10$，$AB=16$，点 $D$ 是边 $AB$ 上的一动点（且 $AD>BD$），若 $\triangle BCD$ 是“差直角三角形”，则 $CD$ 的长为
  \end{minipage}\hfill
  \begin{diagramcoltikz}{65mm}{}
  \includegraphics[width=\linewidth]{\detokenize{../2026-07-24-上海初三试卷原题库/staging/2026-CHONGMING-TERM/items/Q018/assets/manual-prompt-1-e17b2ac1b6da10e0.png}}
  \end{diagramcoltikz}
  \par\medskip
  \answerarea[48mm]
\end{question}

\newpage

% ---------- 第 15 题 (青浦·选择6) ----------
\needspace{10\baselineskip}

\begin{question}[points=4]
    \hfill{\color{edu-blue!60}\small\textit{青浦·选择6}}\par
    学校在操场上举行升旗仪式，每个学生均站立在旗杆前方的水平地面上面向国旗行注目礼。已知国旗的初始位置距地面高度均大于或等于每个学生的身高。如果将看向国旗的视线与水平视线所形成的夹角定义为注视角，那么国旗从初始位置开始匀速上升至顶端的过程中，以下说法正确的是（　）。
    \par
    \begin{choices}[columns=1]
      \item 每个学生的注视角大小不变
      \item 每个学生的注视角逐渐减小
      \item 每个学生的注视角逐渐增大
      \item 同一时刻相同身高学生的注视角相等
    \end{choices}
  \par\medskip
  \answerarea[40mm]
\end{question}

% ---------- 第 16 题 (青浦·填空18) ----------
\needspace{10\baselineskip}

\begin{question}[points=4]
\noindent
  \begin{minipage}[t]{\dimexpr\linewidth-46mm-6mm\relax}
    \vspace{0pt}
    \hfill{\color{edu-blue!60}\small\textit{青浦·填空18}}\par
    如图4，已知 $\triangle ABC$ 中，$\tan B=\dfrac{\sqrt2}{3}，\tan C=\sqrt2，BC=4\sqrt2$。点 $M$ 是 $BC$ 中点，点 $D$ 在边 $AB$ 上，联结 $DM$，将 $\triangle BDM$ 沿直线 $DM$ 翻折，点 $B$ 的对应点为点 $E$。联结 $AE$，如果 $AE\parallel BC$，那么 $\angle BMD$ 的度数为 $\fillin$。
  \end{minipage}\hfill
  \begin{diagramcoltikz}{46mm}{}
  \includegraphics[width=\linewidth]{\detokenize{../2026-07-24-上海初三试卷原题库/staging/2026-QINGPU-TERM/items/Q018/assets/prompt-01.png}}
  \end{diagramcoltikz}
  \par\medskip
  \answerarea[48mm]
\end{question}

\newpage

% ---------- 第 17 题 (虹口·选择6) ----------
\needspace{10\baselineskip}

\begin{question}[points=4]
\noindent
  \begin{minipage}[t]{\dimexpr\linewidth-65mm-6mm\relax}
    \vspace{0pt}
    \hfill{\color{edu-blue!60}\small\textit{虹口·选择6}}\par
    如图3，在 $\triangle ABC$ 中，$AB=2AC$，点 $D$ 在边 $AB$ 上，且 $\angle ACD=\angle B$，如果 $AD=2$，那么 $BD$ 的长是
    \par
    \begin{choices}[columns=1]
      \item $3$
      \item $4$
      \item $5$
      \item $6$
    \end{choices}
  \end{minipage}\hfill
  \begin{diagramcoltikz}{65mm}{}
  \includegraphics[width=\linewidth]{\detokenize{../2026-07-24-上海初三试卷原题库/staging/2026-HONGKOU-TERM/items/Q006/assets/prompt-figure-01.png}}
  \end{diagramcoltikz}
  \par\medskip
  \answerarea[40mm]
\end{question}

% ---------- 第 18 题 (虹口·填空18) ----------
\needspace{10\baselineskip}

\begin{question}[points=4]
\noindent
  \begin{minipage}[t]{\dimexpr\linewidth-65mm-6mm\relax}
    \vspace{0pt}
    \hfill{\color{edu-blue!60}\small\textit{虹口·填空18}}\par
    如图8，在 $\triangle ABC$ 中，$BC=5,\cot A=2,\sin B=\frac35$，$D$ 是 $AB$ 的中点。$E$ 是线段 $AC$ 延长线上一点，联结 $BE$，如果四边形 $BDCE$ 的一组对角相等且另一组对角不相等，那么 $AE$ 的长是\fillin。
  \end{minipage}\hfill
  \begin{diagramcoltikz}{65mm}{}
  \includegraphics[width=\linewidth]{\detokenize{../2026-07-24-上海初三试卷原题库/staging/2026-HONGKOU-TERM/items/Q018/assets/prompt-figure-01.png}}
  \end{diagramcoltikz}
  \par\medskip
  \answerarea[48mm]
\end{question}

\newpage

% ---------- 第 19 题 (杨浦·选择6) ----------
\needspace{10\baselineskip}

\begin{question}[points=4]
\noindent
  \begin{minipage}[t]{\dimexpr\linewidth-65mm-6mm\relax}
    \vspace{0pt}
    \hfill{\color{edu-blue!60}\small\textit{杨浦·选择6}}\par
    如图，已知在 $\triangle ABC$ 中，点 $D、E$ 分别在边 $AB、AC$ 上，$DE\parallel BC$，$\angle ACD=\angle B$，下列结论中错误的是
    \par
    \begin{choices}[columns=1]
      \item $\dfrac{DC^2}{BC^2}=\dfrac{AE}{AC}$
      \item $\dfrac{DE^2}{DC^2}=\dfrac{AD}{AB}$
      \item $\dfrac{AD^2}{AC^2}=\dfrac{AE}{AB}$
      \item $\dfrac{EC^2}{BD^2}=\dfrac{DE}{BC}$
    \end{choices}
  \end{minipage}\hfill
  \begin{diagramcoltikz}{65mm}{}
  \includegraphics[width=\linewidth]{\detokenize{../2026-07-24-上海初三试卷原题库/staging/2026-YANGPU-TERM/items/Q006/assets/prompt-figure.png}}
  \end{diagramcoltikz}
  \par\medskip
  \answerarea[40mm]
\end{question}

% ---------- 第 20 题 (杨浦·填空18) ----------
\needspace{10\baselineskip}

\begin{question}[points=4]
\noindent
  \begin{minipage}[t]{\dimexpr\linewidth-65mm-6mm\relax}
    \vspace{0pt}
    \hfill{\color{edu-blue!60}\small\textit{杨浦·填空18}}\par
    如图，平行四边形 $ABCD$ 中，$AB=6$，$BC=5$。过点 $D$ 作 $DE\perp AB$，垂足为点 $E$；将 $\triangle ADE$ 绕点 $E$ 旋转得 $\triangle FGE$，点 $A$ 的对应点 $F$ 落在边 $AD$ 上，点 $D$ 的对应点 $G$ 落在边 $BC$ 上，连接 $FG$ 与 $DE$ 交于点 $H$，那么 $\dfrac{FH}{GH}=\fillin$。
  \end{minipage}\hfill
  \begin{diagramcoltikz}{65mm}{}
  \includegraphics[width=\linewidth]{\detokenize{../2026-07-24-上海初三试卷原题库/staging/2026-YANGPU-TERM/items/Q018/assets/prompt-figure.png}}
  \end{diagramcoltikz}
  \par\medskip
  \answerarea[48mm]
\end{question}

\newpage

% ---------- 第 21 题 (黄浦·选择6) ----------
\needspace{10\baselineskip}

\begin{question}[points=4]
\noindent
  \begin{minipage}[t]{\dimexpr\linewidth-48mm-6mm\relax}
    \vspace{0pt}
    \hfill{\color{edu-blue!60}\small\textit{黄浦·选择6}}\par
    如图，$D、E$ 是 $\triangle ABC$ 边 $AB、AC$ 上的两点，在下列条件中，能够判定 $DE\parallel BC$ 的是（　）。
    \par
    \begin{choices}[columns=1]
      \item $DB=1，CE=2，AD=3，AE=4$
      \item $DB=1，CE=2，AB=3，AC=4$
      \item $DB=1，CE=2，AD=3，AC=4$
      \item $DB=1，CE=2，AB=3，AE=4$
    \end{choices}
  \end{minipage}\hfill
  \begin{diagramcoltikz}{48mm}{}
  \includegraphics[width=\linewidth]{\detokenize{../2026-07-24-上海初三试卷原题库/staging/2026-HUANGPU-TERM/items/Q006/assets/prompt-01.png}}
  \end{diagramcoltikz}
  \par\medskip
  \answerarea[40mm]
\end{question}

% ---------- 第 22 题 (黄浦·填空16) ----------
\needspace{10\baselineskip}

\begin{question}[points=4]
    \hfill{\color{edu-blue!60}\small\textit{黄浦·填空16}}\par
    对于抛物线 $y=ax^2+bx+c$ 及点 $P$，当过点 $P$ 垂直于抛物线对称轴的直线与该抛物线有两个交点，且两个交点位于点 $P$ 的两侧时，称点 $P$ 为该抛物线的内点。现有抛物线 $L_1:y=-x^2+2x+2$ 和 $L_2:y=\dfrac14x^2+\dfrac34x-\dfrac12$，如果点 $M$ 既是 $L_1$ 的内点，又是 $L_2$ 的内点，那么点 $M$ 的纵坐标 $y_M$ 的取值范围是 $\fillin$。
  \par\medskip
  \answerarea[48mm]
\end{question}

\newpage

% ---------- 第 23 题 (徐汇·选择6) ----------
\needspace{10\baselineskip}

\begin{question}[points=4]
\noindent
  \begin{minipage}[t]{\dimexpr\linewidth-65mm-6mm\relax}
    \vspace{0pt}
    \hfill{\color{edu-blue!60}\small\textit{徐汇·选择6}}\par
    如图，矩形 $ABCD$ 中，点 $E$ 在边 $AB$ 上，点 $F$ 在边 $BC$ 的延长线上，$DF\perp DE$，$DE$ 与对角线 $AC$ 交于点 $I$，$EF$ 与对角线 $AC$ 交于点 $G$，与边 $DC$ 交于点 $H$。则下列结论不一定成立的是
    \par
    \begin{choices}[columns=1]
      \item $\triangle DFH\sim\triangle GAE$
      \item $\triangle ADE\sim\triangle DCA$
      \item $\triangle DAI\sim\triangle GEI$
      \item $\triangle DFH\sim\triangle GCH$
    \end{choices}
  \end{minipage}\hfill
  \begin{diagramcoltikz}{65mm}{}
  \includegraphics[width=\linewidth]{\detokenize{../2026-07-24-上海初三试卷原题库/staging/2026-XUHUI-TERM/items/Q006/assets/prompt-figure.png}}
  \end{diagramcoltikz}
  \par\medskip
  \answerarea[40mm]
\end{question}

% ---------- 第 24 题 (徐汇·填空18) ----------
\needspace{10\baselineskip}

\begin{question}[points=4]
\noindent
  \begin{minipage}[t]{\dimexpr\linewidth-65mm-6mm\relax}
    \vspace{0pt}
    \hfill{\color{edu-blue!60}\small\textit{徐汇·填空18}}\par
    如图，在 $\triangle ABC$ 中，$AB=AC=5$，$\cos C=\frac45$，将 $\triangle ABC$ 绕点 $A$ 逆时针旋转得到 $\triangle ADE$，点 $B、C$ 分别与点 $D、E$ 对应，边 $AD、DE$ 分别与原三角形底边 $BC$ 交于点 $F、G$。当 $\triangle DFG$ 是等腰三角形时，$FG$ 的长为
  \end{minipage}\hfill
  \begin{diagramcoltikz}{65mm}{}
  \includegraphics[width=\linewidth]{\detokenize{../2026-07-24-上海初三试卷原题库/staging/2026-XUHUI-TERM/items/Q018/assets/prompt-figure.png}}
  \end{diagramcoltikz}
  \par\medskip
  \answerarea[48mm]
\end{question}

\newpage

% ---------- 第 25 题 (宝山·选择6) ----------
\needspace{10\baselineskip}

\begin{question}[points=4]
\noindent
  \begin{minipage}[t]{\dimexpr\linewidth-65mm-6mm\relax}
    \vspace{0pt}
    \hfill{\color{edu-blue!60}\small\textit{宝山·选择6}}\par
    小孔 $O$ 到人 $AB$ 的距离、小孔 $O$ 到所成像 $CD$ 的距离均为 $6$ 米。要使像 $CD$ 的长度变为原来的 $1.5$ 倍，下列操作正确的是（　）。
    \par
    \begin{choices}[columns=1]
      \item 人向暗室后退 $2$ 米
      \item 人向暗室前进 $2$ 米
      \item 人向暗室后退 $4$ 米
      \item 人向暗室前进 $4$ 米
    \end{choices}
  \end{minipage}\hfill
  \begin{diagramcoltikz}{65mm}{}
  \includegraphics[width=\linewidth]{\detokenize{../2026-07-24-上海初三试卷原题库/staging/2026-BAOSHAN-TERM/items/Q006/assets/prompt-figure.png}}
  \end{diagramcoltikz}
  \par\medskip
  \answerarea[40mm]
\end{question}

% ---------- 第 26 题 (宝山·填空18) ----------
\needspace{10\baselineskip}

\begin{question}[points=4]
    \hfill{\color{edu-blue!60}\small\textit{宝山·填空18}}\par
    定义：有两个内角的差为 $90^\circ$ 的三角形叫做“差直三角形”。在 $\triangle ABC$ 中，$\angle B>90^\circ$，$\tan C=\dfrac34$，$BC=4$。如果 $\triangle ABC$ 是“差直三角形”，那么 $AB$ 的长为 $\fillin$。
  \par\medskip
  \answerarea[48mm]
\end{question}

\newpage

% ---------- 第 27 题 (奉贤·选择6) ----------
\needspace{10\baselineskip}

\begin{question}[points=4]
    \hfill{\color{edu-blue!60}\small\textit{奉贤·选择6}}\par
    一个三角形框架模型的边长分别为 $20$ cm、$30$ cm、$40$ cm。工人师傅要利用两根 $30$ cm 和 $50$ cm 的铁丝做一个与模型相似的三角形，要求以其中一根铁丝为一边，另一根上截出两段（允许有余料）作为另外两边，那么工人师傅做的这个三角形的周长是
    \par
    \begin{choices}[columns=1]
      \item $80$ cm
      \item $67.5$ cm
      \item $60$ cm
      \item $52.5$ cm
    \end{choices}
  \par\medskip
  \answerarea[40mm]
\end{question}

% ---------- 第 28 题 (奉贤·填空18) ----------
\needspace{10\baselineskip}

\begin{question}[points=4]
\noindent
  \begin{minipage}[t]{\dimexpr\linewidth-65mm-6mm\relax}
    \vspace{0pt}
    \hfill{\color{edu-blue!60}\small\textit{奉贤·填空18}}\par
    如图，在 $\triangle ABC$ 中，$AB=AC=5$，$\cos B=\frac35$。将 $\triangle ABC$ 绕点 $A$ 旋转，点 $B$ 的对应点为点 $D$，联结 $AD$ 交边 $BC$ 于点 $E$。如果 $AE=3ED$，那么 $CE$ 的长为 $\fillin$。
  \end{minipage}\hfill
  \begin{diagramcoltikz}{65mm}{}
  \includegraphics[width=\linewidth]{\detokenize{../2026-07-24-上海初三试卷原题库/staging/2026-FENGXIAN-TERM/items/Q018/assets/manual-prompt-1-2cd9e6da1a40cbc4.png}}
  \end{diagramcoltikz}
  \par\medskip
  \answerarea[48mm]
\end{question}

\newpage

% ---------- 第 29 题 (期末练习卷·选择6) ----------
\needspace{10\baselineskip}

\begin{question}[points=4]
\noindent
  \begin{minipage}[t]{\dimexpr\linewidth-0.30\linewidth-6mm\relax}
    \vspace{0pt}
    \hfill{\color{edu-blue!60}\small\textit{期末练习卷·选择6}}\par
    如图3，已知抛物线 $y=ax^2+bx+c\,(a\ne0)$ 的对称轴是直线 $x=1$，且过点 $(-1,0)$，顶点在第一象限，其部分图像如图所示。给出以下结论：① $ab<0$；② $4a+2b+c>0$；③ $3a+c>0$。其中正确的选项是（　）
    \par
    \begin{choices}[columns=1]
      \item ①②
      \item ①③
      \item ②③
      \item ①②③
    \end{choices}
  \end{minipage}\hfill
  \begin{diagramcoltikz}{0.30\linewidth}{}
  \includegraphics[width=\linewidth]{\detokenize{../2026-07-24-上海初三试卷原题库/staging/GEN-TERM/items/Q006/assets/q06-prompt.png}}
  \end{diagramcoltikz}
  \par\medskip
  \answerarea[40mm]
\end{question}

% ---------- 第 30 题 (期末练习卷·填空18) ----------
\needspace{10\baselineskip}

\begin{question}[points=4]
\noindent
  \begin{minipage}[t]{\dimexpr\linewidth-58mm-6mm\relax}
    \vspace{0pt}
    \hfill{\color{edu-blue!60}\small\textit{期末练习卷·填空18}}\par
    如图，在平行四边形 $ABCD$ 中，$AB=6$，$BC=8$，$\angle B=60^\circ$。将平行四边形 $ABCD$ 绕边 $BC$ 的中点 $O$ 旋转到平行四边形 $A'B'C'D'$ 的位置。当点 $B'$ 落在边 $AB$ 上时，边 $A'D'$ 与边 $CD$ 相交于点 $G$，则 $\cot\angle CGB'=$ \underline{\hspace{2cm}}。
  \end{minipage}\hfill
  \begin{diagramcoltikz}{58mm}{}
  \includegraphics[width=\linewidth]{\detokenize{../2026-07-24-上海初三试卷原题库/staging/GEN-TERM/items/Q018/assets/prompt-01.png}}
  \end{diagramcoltikz}
  \par\medskip
  \answerarea[48mm]
\end{question}

\end{document}